%% secoes/introducao/03-malformacao.tex — incluído por main.tex

%\subsection{Malformação congênita}\label{subsec:intro-malf}

As malformações congênitas, ou anomalias congênitas, são alterações estruturais ou funcionais que se originam durante a vida intrauterina e podem ser identificadas antes do nascimento, ao nascer ou mais tarde na infância \cite{who2023congenital}. Estima-se que cerca de 240 mil recém-nascidos morram por ano no mundo, nos primeiros 28 dias de vida, em decorrência delas, além de outras 170 mil mortes entre um mês e cinco anos de idade \cite{who2023congenital}. A maioria não tem causa conhecida, e as anomalias graves concentram-se nos países de baixa e média renda \cite{who2023congenital}.
No Brasil, a presença de anomalia congênita ao nascer foi de 8,0 por mil nascidos vivos, valor que os autores consideram afetado pela detecção pré-natal limitada e pela subnotificação \cite{melkan2025}.

%No Brasil, a presença de anomalia congênita é registrada na declaração de nascido vivo. Entre 2011 e 2020, a prevalência ao nascer foi de 8,0 por mil nascidos vivos, valor que os autores consideram afetado pela detecção pré-natal limitada e pela subnotificação \cite{melkan2025}. Ainda que fatores genéticos e biológicos tenham papel central, há associação com condições sociais: em âmbito nacional, a chance de anomalia congênita foi maior entre filhos de mães pretas, sem consultas de pré-natal, com mais de 40 anos e em gestações múltiplas \cite{fernandes2025}. Na epidemia de Zika, os casos da síndrome congênita concentraram-se entre as mulheres socialmente mais vulneráveis \cite{paixao2022}.
